%==================================================================
% KATA PENGANTAR
%==================================================================
\chapter*{KATA PENGANTAR}
\frontmatterentry{KATA PENGANTAR}

\vspace{0.5\baselineskip}

\noindent Puji syukur penulis panjatkan ke hadirat Tuhan Yang Maha Esa
atas rahmat dan karunia-Nya sehingga penulis dapat menyelesaikan
\jenisjudul\ yang berjudul ``\textit{\judulid}''. \jenisjudul\ ini
disusun sebagai salah satu syarat kelulusan Program Studi \jenjang\
\prodi\ \institusi.

Penulis menyampaikan terima kasih kepada:

\begin{packed_enum}
  \item Bapak \ldots, selaku Direktur \institusi;
  \item Bapak \ldots, selaku Ketua Jurusan \jurusan;
  \item Bapak/Ibu \ldots, selaku Ketua Program Studi \prodi;
  \item Bapak/Ibu \pembimbingA\ selaku Dosen Pembimbing I dan
        Bapak/Ibu \pembimbingB\ selaku Dosen Pembimbing II;
  \item Bapak/Ibu dosen dan staf \jurusan\ yang telah memberikan ilmu
        dan bantuan selama proses penyusunan \jenisjudul;
  \item Orang tua dan keluarga penulis yang senantiasa mendoakan dan
        memberi dukungan;
  \item Rekan-rekan mahasiswa yang telah membantu pelaksanaan
        penelitian ini.
\end{packed_enum}

Penulis menyadari bahwa \jenisjudul\ ini masih memiliki kekurangan.
Oleh karena itu, kritik dan saran yang membangun sangat penulis
harapkan. Semoga \jenisjudul\ ini bermanfaat bagi pembaca dan menjadi
rujukan bagi pengembangan penelitian selanjutnya.

\vspace{2\baselineskip}

\begin{flushright}
\begin{tabular}{c}
\kota, \tanggalUjian \\
Penulis \\[2.2cm]
\makebox[5.4cm]{\dotfill} \\
\penulis \\
NIM. \nim \\
\end{tabular}
\end{flushright}
